% =============================================================================
% as_qspi_spec.tex  –  Bit-genaue Spezifikation
%   as_qspi Kernel  /  as_slave_bpi  /  as_qspi_top
%   Infineon / AS-Style
% =============================================================================

\subsection{QSPI Peripheral Specification}

\subsubsection{Authors}
\begin{table}[H]
\caption{Authors}
\label{tab:asqspispec_auth}
\centering
\begin{tabularx}{\textwidth}{|X |X |X |}
  \hline
  Author & Department & From \\
  \hline\hline
  Andreas Siggelkow & Fac. E & April 2026 \\
  \hline
  && \\
  \hline
\end{tabularx}
\end{table}

\begin{table}[H]
\caption{History}
\label{tab:asqspispec_hist}
\centering
\begin{tabularx}{\textwidth}{|X |X |X |}
  \hline
  Date & Name & Content Changes \\
  \hline
  \multicolumn{3}{|l|}{Version of this document: V0.1} \\
  \hline
  2026-04-12 & A.~Siggelkow & First draft (AS, Claude) \\
  \hline
\end{tabularx}
\end{table}

% =============================================================================
\subsection{Read-Write Access Types}
% =============================================================================
\begin{itemize}
  \item \textbf{r}  -- read-only; writes have no effect
  \item \textbf{w}  -- write-only; reads return 0
  \item \textbf{rw} -- read/write; value is retained after reset
  \item \textbf{rh} -- read-only; value is set by hardware
  \item \textbf{X0} -- read always returns 0; write-only access
\end{itemize}

% =============================================================================
\subsection{Module: \texttt{as\_qspi} -- QSPI Peripheral Kernel}
\label{subsec:asqspi_kernel}
% =============================================================================

\subsubsection{Functional Description}

The \texttt{as\_qspi} kernel implements the SPI/QSPI state machine and clock
generator. It is a fully state-driven Moore FSM. No FSM state depends on
interrupt signals; all external behavior is driven by register state only.

The FSM traverses the states IDLE\,$\to$\, CMD\,$\to$\, ADDR\,$\to$\, DUM\,$\to$\, DAT\,%
$\to$\, DONE\,$\to$\, IDLE for a normal transaction. In XIP mode the CMD phase
is omitted after the first transaction.

The SCK clock is generated internally from \texttt{clk\_i} via a symmetric
counter. The SCK frequency is:
\[
  f_{SCK} = \frac{f_{clk}}{2 \times (\text{CLKDIV} + 1)}
\]

The kernel drives \texttt{data\_io[3:0]} as output during CMD, ADDR, and TX-DAT
phases and samples \texttt{data\_io[0]} (single), \texttt{data\_io[1:0]} (dual),
or \texttt{data\_io[3:0]} (quad) during RX-DAT phases.

\paragraph{Bits per Cycle (bpc):}
\begin{itemize}
  \item QUAD=1: bpc = 4
  \item DUAL=1, QUAD=0: bpc = 2
  \item DUAL=0, QUAD=0: bpc = 1 (default)
\end{itemize}

\paragraph{Start Semantics:}
\[
  start\_s = (start\_i == 1) \;\land\; (stat\_busy\_o == 0)
\]
The \texttt{start\_i} signal must be a one-cycle pulse, generated by the BPI/Top
from the CTRL.START write strobe.

\paragraph{XIP Mode:}
When \texttt{ctrl\_reg\_i.xip = 1} the CMD phase is skipped on subsequent
transactions. The controller sets \texttt{xip\_active\_o} after the first
DONE state. XIP is cleared by writing CTRL.XIP\,=\,0; the current transfer
completes before IDLE.

\paragraph{Error Latching:}
\begin{itemize}
  \item TX underflow: \texttt{tx\_empty\_i = 1} while in DAT phase and
    \texttt{tx\_flag = 1} (TX direction)
  \item RX overflow: \texttt{rx\_full\_i = 1} while in DAT phase and
    \texttt{tx\_flag = 0} (RX direction)
  \item Timeout: internal counter reaches zero while kernel is active
\end{itemize}

\subsubsection{I/O Ports -- \texttt{as\_qspi}}

\begin{table}[H]
\caption{\texttt{as\_qspi} Ports -- System and Configuration}
\label{tab:asqspi_ports1}
\centering
\begin{tabularx}{\textwidth}{|l|l|l|X|}
  \hline
  Port & Dir. & Width & Description \\
  \hline\hline
  \texttt{clk\_i}           & in  & 1  & Kernel clock, active rising edge \\
  \hline
  \texttt{rst\_i}           & in  & 1  & Synchronous reset, active high \\
  \hline
  \texttt{start\_i}         & in  & 1  & Start pulse (one cycle wide). Initiates FSM if not busy. \\
  \hline
  \texttt{ctrl\_reg\_i}     & in  & 8  & Packed \texttt{qspi\_ctrl\_t}: \{addr\_len, cpol, cpha, dual, cs\_hold, xip, ddr, quad\} \\
  \hline
  \texttt{cmd\_reg\_i}      & in  & 8  & SPI opcode byte \\
  \hline
  \texttt{addr\_reg\_i}     & in  & 32 & Flash target address \\
  \hline
  \texttt{len\_reg\_i}      & in  & 16 & Transfer length in bytes \\
  \hline
  \texttt{dummy\_reg\_i}    & in  & 6  & Number of dummy cycles \\
  \hline
  \texttt{clkdiv\_reg\_i}   & in  & 8  & SCK clock divider \\
  \hline
  \texttt{timeout\_reg\_i}  & in  & 32 & Timeout threshold in clock cycles; 0 = disabled \\
  \hline
  \texttt{xip\_mode\_bits\_i}& in & 8  & XIP mode byte, driven on IO[3:0] during XIP address phase \\
  \hline
\end{tabularx}
\end{table}

\begin{table}[H]
\caption{\texttt{as\_qspi} Ports -- Status, FIFO and SPI}
\label{tab:asqspi_ports2}
\centering
\begin{tabularx}{\textwidth}{|l|l|l|X|}
  \hline
  Port & Dir. & Width & Description \\
  \hline\hline
  \texttt{stat\_busy\_o}    & out & 1  & FSM is active (not IDLE, not DONE). Level signal. \\
  \hline
  \texttt{stat\_done\_o}    & out & 1  & FSM is in DONE state. High for exactly one clock cycle. \\
  \hline
  \texttt{stat\_error\_o}   & out & 1  & Error latched. Remains until ICR clear. \\
  \hline
  \texttt{stat\_timeout\_o} & out & 1  & Timeout occurred. Remains until ICR clear. \\
  \hline
  \texttt{xip\_active\_o}   & out & 1  & XIP mode is active. \\
  \hline
  \texttt{tx\_empty\_i}     & in  & 1  & TX FIFO empty indication (from TX FIFO or BPI) \\
  \hline
  \texttt{tx\_rd\_o}        & out & 1  & TX FIFO read strobe. One cycle pulse on last TX drive. \\
  \hline
  \texttt{tx\_data\_i}      & in  & 64 & TX data word from TX FIFO \\
  \hline
  \texttt{rx\_full\_i}      & in  & 1  & RX FIFO full indication \\
  \hline
  \texttt{rx\_wr\_o}        & out & 1  & RX FIFO write strobe. One cycle pulse on last RX sample. \\
  \hline
  \texttt{rx\_data\_o}      & out & 64 & RX shift register output. Valid after \texttt{rx\_wr\_o}. \\
  \hline
  \texttt{sck\_o}           & out & 1  & SPI clock output \\
  \hline
  \texttt{cs\_o}            & out & 1  & Chip select (active high). Deasserted in IDLE and DONE. \\
  \hline
  \texttt{data\_io}         & inout & 4 & SPI data bus (tri-state). Driven in CMD/ADDR/TX-DAT, sampled in RX-DAT. \\
  \hline
\end{tabularx}
\end{table}

\subsubsection{Kernel State Machine}

\begin{table}[H]
\caption{\texttt{as\_qspi} FSM States}
\label{tab:asqspi_fsm}
\centering
\begin{tabularx}{\textwidth}{|l|l|l|X|}
  \hline
  State & \texttt{cs\_o} & \texttt{sck\_o} & Condition to leave \\
  \hline\hline
  IDLE\_ST  & 0 & 0 & \texttt{start\_i = 1} \\
  \hline
  CMD\_ST   & 1 & running & \texttt{cnt\_cmd\_r} has reached CMD\_MAX (7); exit on \texttt{sck\_drive\_s} \\
  \hline
  ADDR\_ST  & 1 & running & \texttt{cnt\_addr\_r} has reached \texttt{addr\_max\_s}; exit on \texttt{sck\_drive\_s} \\
  \hline
  DUM\_ST   & 1 & running & \texttt{cnt\_dum\_r + 1 $\geq$ dummy\_reg\_i}; exit on \texttt{sck\_rise\_s} \\
  \hline
  DAT\_ST   & 1 & running & \texttt{cnt\_dat\_r + bpc\_s $>$ dat\_max\_s}; exit on last sample (RX) or last drive (TX) \\
  \hline
  DONE\_ST  & 0 & 0 & Unconditional, one cycle \\
  \hline
\end{tabularx}
\end{table}

\begin{table}[H]
\caption{\texttt{as\_qspi} Internal Counters}
\label{tab:asqspi_cnt}
\centering
\begin{tabularx}{\textwidth}{|l|l|X|}
  \hline
  Counter & Width & Description \\
  \hline\hline
  \texttt{sck\_cnt\_r}  & 8 & SCK clock divider counter. Resets to \texttt{clkdiv\_reg\_i}. \\
  \hline
  \texttt{cnt\_cmd\_r}  & 3 & CMD phase bit counter. Counts 0..7 (always 8 bits, single wire). \\
  \hline
  \texttt{cnt\_addr\_r} & 5 & ADDR phase bit counter. Max = 23 (3-byte) or 31 (4-byte). \\
  \hline
  \texttt{cnt\_dum\_r}  & 6 & DUMMY cycle counter. Counts up to \texttt{dummy\_reg\_i}. \\
  \hline
  \texttt{cnt\_dat\_r}  & 20 & DAT phase bit counter. Counts in steps of \texttt{bpc\_s}. Max = \texttt{len\_reg\_i*8 - 1}. \\
  \hline
\end{tabularx}
\end{table}

\subsubsection{Timing Constraints}

\begin{itemize}
  \item All configuration registers (\texttt{cmd\_reg\_i}, \texttt{addr\_reg\_i},
    \texttt{len\_reg\_i}, etc.) must be stable before \texttt{start\_i} is asserted
    and must remain stable for the duration of the transfer (BUSY=1).
  \item \texttt{clkdiv\_reg\_i} must not change while BUSY=1.
  \item \texttt{tx\_data\_i} must be valid when \texttt{tx\_rd\_o} fires. The BPI
    preloads the TX shift register one cycle before DAT\_ST.
  \item \texttt{rx\_data\_o} is valid two clock cycles after \texttt{rx\_wr\_o}
    (NBA settling time in simulation; in synthesis: same cycle is valid).
\end{itemize}

% =============================================================================
\subsection{Module: \texttt{as\_slave\_bpi} -- Wishbone Slave Bus Peripheral Interface}
\label{subsec:asqspi_bpi}
% =============================================================================

\subsubsection{Functional Description}

The \texttt{as\_slave\_bpi} module implements a Wishbone Classic Compliant Slave
BPI. It provides zero-wait-state single read/write transfers to the attached
peripheral (QSPI top level). The BPI is fully combinatorial on the bus side;
ACK is asserted in the same clock cycle as STB.

\paragraph{Protocol:}
A valid Wishbone transfer requires \texttt{CYC\_I = 1} and \texttt{STB\_I = 1}
simultaneously. All SEL bits must be asserted (\texttt{SEL\_I = 0xFF}) for a
transfer to be acknowledged. The ACK is combinatorial:
\[
  ACK\_O = CYC\_I \;\land\; STB\_I \;\land\; (\&SEL\_I)
\]

\paragraph{Write:}
When \texttt{WE\_I = 1}, \texttt{WR\_O} is asserted to the peripheral for the
duration of the bus cycle. The address decoder in the Top level routes the write
to the appropriate register.

\paragraph{Read:}
When \texttt{WE\_I = 0}, \texttt{RD\_O} is asserted. The peripheral must provide
read data combinatorially on \texttt{DAT\_FROM\_CORE\_I}; this is passed
directly to \texttt{WB\_S\_DAT\_O}.

\subsubsection{I/O Ports -- \texttt{as\_slave\_bpi}}

\begin{table}[H]
\caption{\texttt{as\_slave\_bpi} Ports}
\label{tab:asbpi_ports}
\centering
\begin{tabularx}{\textwidth}{|l|l|l|X|}
  \hline
  Port & Dir. & Width & Description \\
  \hline\hline
  \texttt{rst\_i}             & in  & 1  & System reset, active high (unused internally; for consistency) \\
  \hline
  \texttt{clk\_i}             & in  & 1  & System clock (unused internally; for consistency) \\
  \hline
  \multicolumn{4}{|l|}{\textit{Peripheral (kernel) side:}} \\
  \hline
  \texttt{addr\_o}            & out & 64 & Decoded byte offset address, directly from \texttt{wb\_s\_addr\_i} \\
  \hline
  \texttt{dat\_from\_core\_i} & in  & 64 & Read data from peripheral, presented combinatorially \\
  \hline
  \texttt{dat\_to\_core\_o}   & out & 64 & Write data to peripheral, directly from \texttt{wb\_s\_dat\_i} \\
  \hline
  \texttt{wr\_o}              & out & 1  & Write enable to peripheral. Combinatorial. \\
  \hline
  \texttt{rd\_o}              & out & 1  & Read enable to peripheral. Combinatorial. \\
  \hline
  \multicolumn{4}{|l|}{\textit{Wishbone slave side:}} \\
  \hline
  \texttt{wb\_s\_addr\_i}     & in  & 64 & Wishbone address (byte offset from peripheral base) \\
  \hline
  \texttt{wb\_s\_dat\_i}      & in  & 64 & Wishbone write data \\
  \hline
  \texttt{wb\_s\_dat\_o}      & out & 64 & Wishbone read data \\
  \hline
  \texttt{wb\_s\_we\_i}       & in  & 1  & Wishbone write enable \\
  \hline
  \texttt{wb\_s\_sel\_i}      & in  & 8  & Wishbone byte select (all bytes must be selected) \\
  \hline
  \texttt{wb\_s\_stb\_i}      & in  & 1  & Wishbone strobe (valid cycle) \\
  \hline
  \texttt{wb\_s\_ack\_o}      & out & 1  & Wishbone acknowledge (combinatorial, zero wait state) \\
  \hline
  \texttt{wb\_s\_cyc\_i}      & in  & 1  & Wishbone cycle (high for complete bus transaction) \\
  \hline
\end{tabularx}
\end{table}

\subsubsection{BPI Transfer Timing}

\begin{table}[H]
\caption{Wishbone Classic Zero-Wait-State Timing}
\label{tab:asbpi_timing}
\centering
\begin{tabularx}{\textwidth}{|l|X|}
  \hline
  Cycle & Description \\
  \hline\hline
  T0 & Master asserts CYC, STB, ADR, DAT (write) or ADR (read), WE, SEL \\
  \hline
  T0 & Slave combinatorially asserts ACK. Read data appears on DAT\_O. \\
  \hline
  T1 & Master samples ACK and DAT\_O on rising clock edge. \\
  \hline
  T1 & Master deasserts STB (and CYC for single transfers). \\
  \hline
\end{tabularx}
\end{table}

\textbf{Note:} The peripheral register file is clocked (\texttt{always\_ff}),
so a write takes effect on the posedge of the clock during which \texttt{WR\_O}
is asserted. Read data must be available combinatorially before the posedge.

% =============================================================================
\subsection{Module: \texttt{as\_qspi\_top} -- QSPI Top Level}
\label{subsec:asqspi_top}
% =============================================================================

\subsubsection{Functional Description}

\texttt{as\_qspi\_top} integrates the Wishbone BPI (\texttt{as\_slave\_bpi}),
the QSPI kernel (\texttt{as\_qspi}), the TX FIFO, and the RX FIFO into a
complete memory-mapped QSPI peripheral. It decodes the register offset address,
routes read/write strobes to the appropriate registers, and implements the
interrupt logic.

\paragraph{Register Address Decoding:}
The BPI passes the raw bus address directly as a byte offset. The Top level
compares this offset against the constants defined in Table~\ref{tab:asqspitop_regmap}
using combinatorial equality. There is no base-address stripping; the bus master
must provide the byte offset from the peripheral's base address.

\paragraph{Start Pulse Generation:}
Writing \texttt{CTRL.START (bit 3) = 1} causes the Top level to register the
write strobe for one clock cycle:
\[
  start\_r \;\leftarrow\; wr\_ctrl \;\land\; data\_wr\_s[3]
\]
\[
  start\_s = start\_r \quad \text{(one-cycle pulse, registered)}
\]

\paragraph{FIFO Flush:}
\texttt{CTRL.TX\_FLUSH (bit 1)} and \texttt{CTRL.RX\_FLUSH (bit 2)} are
combinatorial and active only during the write cycle:
\[
  tx\_flush\_s = wr\_ctrl \;\land\; data\_wr\_s[1]
\]
\[
  rx\_flush\_s = wr\_ctrl \;\land\; data\_wr\_s[2]
\]
Flushing while BUSY is asserted is implementation-defined; in this
implementation the flush takes effect immediately.

\paragraph{RX FIFO Write Pipeline:}
The kernel outputs \texttt{rx\_wr\_o} and \texttt{rx\_data\_o} (\texttt{rx\_shift\_r})
as NBA assignments in the same clock cycle. To avoid a simulation race condition
(and to provide a registered interface for synthesis), a two-cycle pipeline
is used:
\begin{itemize}
  \item Cycle N: \texttt{rx\_wr\_kernel\_s = 1}, \texttt{rx\_shift\_r} NBA pending.
  \item Cycle N+1: \texttt{rx\_wr\_d1 = 1}, \texttt{rx\_data\_kernel\_s} captured
    into \texttt{rx\_data\_snap} (shift register now settled).
  \item Cycle N+2: \texttt{rx\_wr\_d2 = 1}, \texttt{rx\_data\_snap} written to RX FIFO.
\end{itemize}

\paragraph{RX FIFO Pop (read-before-pop):}
Reading \texttt{RXDATA} presents \texttt{mem[rd\_ptr\_r]} combinatorially and
schedules the pointer advance one cycle later via \texttt{rx\_pop\_r}. This
separates data presentation from pointer advancement and is required for
correct operation with the asynchronous-read \texttt{as\_fifo}.

\subsubsection{I/O Ports -- \texttt{as\_qspi\_top}}

\begin{table}[H]
\caption{\texttt{as\_qspi\_top} Ports}
\label{tab:asqspitop_ports}
\centering
\begin{tabularx}{\textwidth}{|l|l|l|X|}
  \hline
  Port & Dir. & Width & Description \\
  \hline\hline
  \texttt{rst\_i}       & in  & 1  & Synchronous reset, active high \\
  \hline
  \texttt{clk\_i}       & in  & 1  & System and kernel clock \\
  \hline
  \texttt{wbdAddr\_i}   & in  & 64 & Wishbone address (byte offset from QSPI base) \\
  \hline
  \texttt{wbdDat\_i}    & in  & 64 & Wishbone write data \\
  \hline
  \texttt{wbdDat\_o}    & out & 64 & Wishbone read data \\
  \hline
  \texttt{wbdWe\_i}     & in  & 1  & Wishbone write enable \\
  \hline
  \texttt{wbdSel\_i}    & in  & 8  & Wishbone byte select \\
  \hline
  \texttt{wbdStb\_i}    & in  & 1  & Wishbone strobe \\
  \hline
  \texttt{wbdAck\_o}    & out & 1  & Wishbone acknowledge \\
  \hline
  \texttt{wbdCyc\_i}    & in  & 1  & Wishbone cycle \\
  \hline
  \texttt{sck\_o}       & out & 1  & SPI clock to NOR-Flash \\
  \hline
  \texttt{cs\_o}        & out & 1  & Chip select to NOR-Flash (active high) \\
  \hline
  \texttt{data\_io}     & inout & 4 & QSPI data bus to/from NOR-Flash \\
  \hline
  \texttt{qspi\_irq\_o} & out & 1  & Interrupt request to CPU interrupt controller \\
  \hline
\end{tabularx}
\end{table}

\subsubsection{Register Map}

\begin{table}[H]
\caption{\texttt{as\_qspi\_top} Register Map}
\label{tab:asqspitop_regmap}
\centering
\begin{tabularx}{\textwidth}{|l|l| X| l| X|}
  \hline
  Offset & Symbol & Reset & Access & Description \\
  \hline\hline
  0x00 (0)   & ID        & 0x00000000 \_00000010 & r   & Peripheral identification register \\
  \hline
  0x08 (8)   & CTRL      & 0x00000000 \_00000000 & rw  & Control register (START, FLUSH, mode bits) \\
  \hline
  0x10 (16)  & CMD       & 0x00000000 \_00000000 & rw  & SPI opcode \\
  \hline
  0x18 (24)  & ADDR      & 0x00000000 \_00000000 & rw  & Flash target address \\
  \hline
  0x20 (32)  & LEN       & 0x00000000 \_00000000 & rw  & Transfer length in bytes \\
  \hline
  0x28 (40)  & DUMMY     & 0x00000000 \_00000000 & rw  & Dummy cycles \\
  \hline
  0x30 (48)  & CLKDIV    & 0x00000000 \_00000000 & rw  & SCK clock divider \\
  \hline
  0x38 (56)  & TIMEOUT   & 0x00000000 \_00000000 & rw  & Timeout threshold \\
  \hline
  0x40 (64)  & ISR       & 0x00000000 \_00000000 & w   & Interrupt Set Register (w1 $\to$ set RIS) \\
  \hline
  0x48 (72)  & RIS       & 0x00000000 \_00000000 & rh  & Raw Interrupt Status Register \\
  \hline
  0x50 (80)  & IMSC      & 0xFFFFFFFF \_FFFFFFFF & rw  & Interrupt Mask Control Register \\
  \hline
  0x58 (88)  & MIS       & 0x00000000 \_00000000 & rh  & Masked Interrupt Status (= RIS \& IMSC) \\
  \hline
  0x60 (96)  & ICR       & 0x00000000 \_00000000 & w   & Interrupt Clear Register (w1 $\to$ clear RIS) \\
  \hline
  0x68 (104) & RXDATA    & 0x00000000 \_00000000 & r   & RX FIFO pop (read-before-pop) \\
  \hline
  0x70 (112) & TXDATA    & 0x00000000 \_00000000 & w   & TX FIFO push \\
  \hline
  0x78 (120) & FIFOSTAT  & 0x00000000 \_00000000 & rh  & FIFO fill levels and flags \\
  \hline
  0x80 (128) & XIPMODE   & 0x00000000 \_000000A0 & rw  & XIP mode byte (Winbond default: 0xA0) \\
  \hline
  0x88 (136) & STATUS    & 0x00000000 \_00000000 & rh  & Kernel status flags \\
  \hline
\end{tabularx}
\end{table}

% =============================================================================
\subsection{Register Descriptions}
% =============================================================================

\subsubsection{System Registers}

\paragraph{ID} -:
Identification register of the QSPI peripheral. Read-only; writes have no effect.

\asregkopf{ID register}{$0000000000000010_H$}
\asregfourxsixteen{\multicolumn{16}{|c|}{0}}
                             {\multicolumn{16}{|c|}{r}}
                             {\multicolumn{16}{|c|}{0}}
                             {\multicolumn{16}{|c|}{r}}
                             {\multicolumn{16}{|c|}{0}}
                             {\multicolumn{16}{|c|}{r}}
                             {\multicolumn{8}{|c|}{0}&\multicolumn{8}{c|}{PERIPH\_ID}}
                             {\multicolumn{8}{|c|}{r}&\multicolumn{8}{c|}{r}}
\asregdesc{- & 63:8 & r & Reserved, reads as 0 \\
           \hline
           PERIPH\_ID & 7:0 & r & Peripheral identification number. Fixed value \texttt{0x10}.}

% =============================================================================
\subsubsection{Special Function Registers}

\paragraph{CTRL} -:
Main control register. Controls SPI mode, address length, clock polarity/phase,
XIP, DDR, CS-hold, FIFO flush, and transfer start.
All configuration bits must be set before asserting START.

\asregkopf{Control register}{$0000000000000000_H$}
\asregfourxsixteen{\multicolumn{16}{|c|}{0}}
                             {\multicolumn{16}{|c|}{r}}
                             {\multicolumn{16}{|c|}{0}}
                             {\multicolumn{16}{|c|}{r}}
                             {\multicolumn{16}{|c|}{0}}
                             {\multicolumn{16}{|c|}{r}}
                             {0&0&RF&TF&AL&CP&CH&DU&CS&XI&DD&QU&ST&0&0&0}
                             {r&r&w&w&rw&rw&rw&rw&rw&rw&rw&rw&w&r&r&r}
\asregdesc{- & 63:14 & r & Reserved \\
           \hline
           RF & 13 & w  & RX\_FLUSH. Write 1 to flush RX FIFO. Self-clearing. \\
           \hline
           TF & 12 & w  & TX\_FLUSH. Write 1 to flush TX FIFO. Self-clearing. \\
           \hline
           AL & 11 & rw & ADDR\_LEN. 0 = 3-byte address (24\,bit); 1 = 4-byte (32\,bit). \\
           \hline
           CP & 10 & rw & CPOL. Clock polarity. 0 = idle low (Mode\,0); 1 = idle high (Mode\,3). \\
           \hline
           CH &  9 & rw & CPHA. Clock phase. 0 = sample on leading edge; 1 = sample on trailing edge. \\
           \hline
           DU &  8 & rw & DUAL. Enable dual-wire data (1-1-2). Ignored when QUAD=1. \\
           \hline
           CS &  7 & rw & CS\_HOLD. Keep CS asserted between consecutive transactions. \\
           \hline
           XI &  6 & rw & XIP. Enable execute-in-place mode. \\
           \hline
           DD &  5 & rw & DDR. Enable double data rate. \\
           \hline
           QU &  4 & rw & QUAD. Enable quad-wire data (1-1-4 / 1-4-4). Overrides DUAL. \\
           \hline
           ST &  3 & w  & START. Write 1 to initiate a transfer when BUSY=0. Self-clearing. \\
           \hline
           -  & 2:0 & r & Reserved.}

\paragraph{CMD} -:
The SPI/QSPI command opcode to be sent in the CMD phase.
\asregkopf{Command register}{$0000000000000000_H$}
\asregfourxsixteen{\multicolumn{16}{|c|}{0}}
                             {\multicolumn{16}{|c|}{r}}
                             {\multicolumn{16}{|c|}{0}}
                             {\multicolumn{16}{|c|}{r}}
                             {\multicolumn{16}{|c|}{0}}
                             {\multicolumn{16}{|c|}{r}}
                             {0&0&0&0&0&0&0&0&\multicolumn{8}{c|}{OPCODE}}
                             {r&r&r&r&r&r&r&r&\multicolumn{8}{c|}{rw}}
\asregdesc{- & 63:8 & r & Reserved \\
           \hline
           OPCODE & 7:0 & rw & SPI command opcode. Sent MSB-first in single-wire CMD phase.}

\paragraph{ADDR} -:
Target byte address in the NOR-Flash device.
\asregkopf{Address register}{$0000000000000000_H$}
\asregfourxsixteen{\multicolumn{16}{|c|}{0}}
                             {\multicolumn{16}{|c|}{r}}
                             {\multicolumn{16}{|c|}{0}}
                             {\multicolumn{16}{|c|}{r}}
                             {\multicolumn{16}{|c|}{ADDR[31:16]}}
                             {\multicolumn{16}{|c|}{rw}}
                             {\multicolumn{16}{|c|}{ADDR[15:0]}}
                             {\multicolumn{16}{|c|}{rw}}
\asregdesc{- & 63:32 & r & Reserved \\
           \hline
           ADDR & 31:0 & rw & Flash byte address. Sent MSB-first. Number of address bytes determined by CTRL.AL.}

\paragraph{LEN} -:
Transfer data length in bytes.
\asregkopf{Length register}{$0000000000000000_H$}
\asregfourxsixteen{\multicolumn{16}{|c|}{0}}
                             {\multicolumn{16}{|c|}{r}}
                             {\multicolumn{16}{|c|}{0}}
                             {\multicolumn{16}{|c|}{r}}
                             {\multicolumn{16}{|c|}{0}}
                             {\multicolumn{16}{|c|}{r}}
                             {\multicolumn{16}{|c|}{LEN}}
                             {\multicolumn{16}{|c|}{rw}}
\asregdesc{- & 63:16 & r & Reserved \\
           \hline
           LEN & 15:0 & rw & Data transfer length in bytes. Must be $\geq 1$.}

\paragraph{DUMMY} -:
Number of dummy clock cycles between ADDR and DATA phases.
\asregkopf{Dummy register}{$0000000000000000_H$}
\asregfourxsixteen{\multicolumn{16}{|c|}{0}}
                             {\multicolumn{16}{|c|}{r}}
                             {\multicolumn{16}{|c|}{0}}
                             {\multicolumn{16}{|c|}{r}}
                             {\multicolumn{16}{|c|}{0}}
                             {\multicolumn{16}{|c|}{r}}
                             {0&0&0&0&0&0&0&0&0&0&\multicolumn{6}{c|}{DUMMY\_CYC}}
                             {r&r&r&r&r&r&r&r&r&r&\multicolumn{6}{c|}{rw}}
\asregdesc{- & 63:6 & r & Reserved \\
           \hline
           DUMMY\_CYC & 5:0 & rw & Dummy cycle count. 0 = no dummy phase (ADDR directly followed by DATA). Maximum 63.}

\paragraph{CLKDIV} -:
SCK clock divider. Determines the SPI clock frequency.
\[
  f_{SCK} = \frac{f_{clk}}{2 \times (\text{CLKDIV} + 1)}
\]
Must not be changed while BUSY=1.
\asregkopf{Clock divider register}{$0000000000000000_H$}
\asregfourxsixteen{\multicolumn{16}{|c|}{0}}
                             {\multicolumn{16}{|c|}{r}}
                             {\multicolumn{16}{|c|}{0}}
                             {\multicolumn{16}{|c|}{r}}
                             {\multicolumn{16}{|c|}{0}}
                             {\multicolumn{16}{|c|}{r}}
                             {\multicolumn{8}{|c|}{0}&\multicolumn{8}{c|}{CLKDIV}}
                             {\multicolumn{8}{|c|}{r}&\multicolumn{8}{c|}{rw}}
\asregdesc{- & 63:8 & r & Reserved \\
           \hline
           CLKDIV & 7:0 & rw & SCK divider. Reset 0x00 = $f_{clk}/2$ (maximum SCK).}

\paragraph{TIMEOUT} -:
Timeout threshold in kernel clock cycles. When the FSM is active and the
down-counter reaches zero, a timeout error is raised. Set to 0 to disable.
\asregkopf{Timeout register}{$0000000000000000_H$}
\asregfourxsixteen{\multicolumn{16}{|c|}{0}}
                             {\multicolumn{16}{|c|}{r}}
                             {\multicolumn{16}{|c|}{0}}
                             {\multicolumn{16}{|c|}{r}}
                             {\multicolumn{16}{|c|}{TIMEOUT\_THR[31:16]}}
                             {\multicolumn{16}{|c|}{rw}}
                             {\multicolumn{16}{|c|}{TIMEOUT\_THR[15:0]}}
                             {\multicolumn{16}{|c|}{rw}}
\asregdesc{- & 63:32 & r & Reserved \\
           \hline
           TIMEOUT\_THR & 31:0 & rw & Timeout threshold in clock cycles. 0x00000000 = disabled.}

\paragraph{TXDATA} -:
Writing to this register pushes data into the TX FIFO. If the TX FIFO is full,
the write is silently discarded; the Wishbone ACK is still asserted.
Reading returns 0.
\asregkopf{TX data register (write pushes to TX FIFO)}{$0000000000000000_H$}
\asregfourxsixteen{\multicolumn{16}{|c|}{TX\_DATA[63:48]}}
                             {\multicolumn{16}{|c|}{w}}
                             {\multicolumn{16}{|c|}{TX\_DATA[47:32]}}
                             {\multicolumn{16}{|c|}{w}}
                             {\multicolumn{16}{|c|}{TX\_DATA[31:16]}}
                             {\multicolumn{16}{|c|}{w}}
                             {\multicolumn{16}{|c|}{TX\_DATA[15:0]}}
                             {\multicolumn{16}{|c|}{w}}
\asregdesc{TX\_DATA & 63:0 & w & Write data word pushed into TX FIFO. MSB transmitted first. Reads return 0.}

\paragraph{RXDATA} -:
Reading this register pops one 64-bit word from the RX FIFO and returns it.
The pop (pointer advance) is delayed by one clock cycle (read-before-pop
semantics). If the RX FIFO is empty, the read returns 0x0000000000000000 and
no pop occurs. Writing has no effect.
\asregkopf{RX data register (read pops from RX FIFO)}{$0000000000000000_H$}
\asregfourxsixteen{\multicolumn{16}{|c|}{RX\_DATA[63:48]}}
                             {\multicolumn{16}{|c|}{r}}
                             {\multicolumn{16}{|c|}{RX\_DATA[47:32]}}
                             {\multicolumn{16}{|c|}{r}}
                             {\multicolumn{16}{|c|}{RX\_DATA[31:16]}}
                             {\multicolumn{16}{|c|}{r}}
                             {\multicolumn{16}{|c|}{RX\_DATA[15:0]}}
                             {\multicolumn{16}{|c|}{r}}
\asregdesc{RX\_DATA & 63:0 & r & Read data word popped from RX FIFO. MSB received first. Writes have no effect.}

\paragraph{FIFOSTAT} -:
FIFO fill level and status flags. Read-only. Updated combinatorially.
\asregkopf{FIFO status register}{$0000000000000000_H$}
\asregfourxsixteen{\multicolumn{16}{|c|}{0}}
                             {\multicolumn{16}{|c|}{r}}
                             {\multicolumn{16}{|c|}{0}}
                             {\multicolumn{16}{|c|}{r}}
                             {0&0&0&0&0&0&RF&RH&0&0&0&\multicolumn{5}{c|}{RX\_LEV[4:0]}}
                             {r&r&r&r&r&r&rh&rh&r&r&r&\multicolumn{5}{c|}{rh}}
                             {0&0&0&0&0&0&TF&TH&0&0&0&\multicolumn{5}{c|}{TX\_LEV[4:0]}}
                             {r&r&r&r&r&r&rh&rh&r&r&r&\multicolumn{5}{c|}{rh}}
\asregdesc{- & 63:26 & r & Reserved \\
           \hline
           RF & 25 & rh & RX FIFO full \\
           \hline
           RH & 24 & rh & RX FIFO $\geq$ half full (level $\geq$ FIFO\_DEPTH/2) \\
           \hline
           RX\_LEV & 20:16 & rh & RX FIFO fill level (0\ldots16) \\
           \hline
           - & 15:10 & r & Reserved \\
           \hline
           TF & 9 & rh & TX FIFO full \\
           \hline
           TH & 8 & rh & TX FIFO $<$ half full (level $<$ FIFO\_DEPTH/2) \\
           \hline
           TX\_LEV & 4:0 & rh & TX FIFO fill level (0\ldots16)}

\paragraph{XIPMODE} -:
XIP mode byte driven on IO[7:4] (upper nibble on IO[3:0]) during the XIP
address phase to terminate the mode register entry sequence
(e.g.\ Winbond W25Q128JV: 0xA0).
\asregkopf{XIP mode byte register}{$00000000000000A0_H$}
\asregfourxsixteen{\multicolumn{16}{|c|}{0}}
                             {\multicolumn{16}{|c|}{r}}
                             {\multicolumn{16}{|c|}{0}}
                             {\multicolumn{16}{|c|}{r}}
                             {\multicolumn{16}{|c|}{0}}
                             {\multicolumn{16}{|c|}{r}}
                             {\multicolumn{8}{|c|}{0}&\multicolumn{8}{c|}{XIP\_BYTE}}
                             {\multicolumn{8}{|c|}{r}&\multicolumn{8}{c|}{rw}}
\asregdesc{- & 63:8 & r & Reserved \\
           \hline
           XIP\_BYTE & 7:0 & rw & Mode byte for XIP address phase. Default 0xA0 (Winbond W25Q128JV).}

\paragraph{STATUS} -:
Kernel status register. All bits are read-only and updated combinatorially.

\textbf{Note on DONE:} \texttt{STATUS.DONE} reflects the combinatorial
\texttt{stat\_done\_o} signal which is high only while the kernel is in DONE\_ST
(one clock cycle). It is \emph{not} suitable for polling. Use
\texttt{RIS.DO} instead, which latches the DONE event and holds it until
cleared via ICR.
\asregkopf{Status register}{$0000000000000000_H$}
\asregfourxsixteen{\multicolumn{16}{|c|}{0}}
                             {\multicolumn{16}{|c|}{r}}
                             {\multicolumn{16}{|c|}{0}}
                             {\multicolumn{16}{|c|}{r}}
                             {\multicolumn{16}{|c|}{0}}
                             {\multicolumn{16}{|c|}{r}}
                             {0&0&0&0&0&0&0&0&0&0&0&TI&ER&XI&DO&BU}
                             {r&r&r&r&r&r&r&r&r&r&r&rh&rh&rh&rh&rh}
\asregdesc{- & 63:5 & r & Reserved \\
           \hline
           TI & 4 & rh & TIMEOUT. Timeout occurred; latched until ICR.TIMEOUT is written. \\
           \hline
           ER & 3 & rh & ERROR. Error latched; latched until ICR.ERROR is written. \\
           \hline
           XI & 2 & rh & XIP\_ACTIVE. XIP mode is currently active. \\
           \hline
           DO & 1 & rh & DONE. High for exactly one clock cycle when FSM leaves DAT\_ST. For polling use RIS.DO. \\
           \hline
           BU & 0 & rh & BUSY. FSM is active (not IDLE).}

% =============================================================================
\subsubsection{Interrupt Registers}
% =============================================================================

\paragraph{Interrupt Request Flow} -:
Figure~\ref{fig:asqspispec_irq} shows the interrupt flow. A rising edge on
an interrupt source sets the corresponding RIS bit. The ICR clears RIS bits,
the ISR forces them. The IRQ output is the OR-reduction of the masked
status (MIS = RIS \& IMSC).
\[
  MIS = RIS \;\&\; IMSC
\]
\[
  IRQ = |~MIS
\]

\paragraph{Interrupt Sources} -:
\begin{table}[H]
\caption{QSPI Interrupt Sources}
\label{tab:asqspispec_irqsrc}
\centering
\begin{tabularx}{\textwidth}{|l|l|X|X|}
  \hline
  RIS Bit & Name & Source & Latch Semantics \\
  \hline\hline
  0 & DONE     & Rising edge of \texttt{stat\_done\_s} (FSM enters DONE\_ST) & Set on edge; hold until ICR \\
  \hline
  1 & RXHALF   & Rising edge of \texttt{rx\_half\_s} (RX FIFO $\geq$ half full) & Set on edge; hold until ICR \\
  \hline
  2 & TXHALF   & Rising edge of \texttt{tx\_half\_s} (TX FIFO $<$ half full)   & Set on edge; hold until ICR \\
  \hline
  3 & ERROR    & Rising edge of \texttt{stat\_error\_s}                         & Set on edge; hold until ICR \\
  \hline
  4 & TIMEOUT  & Rising edge of \texttt{stat\_timeout\_s}                       & Set on edge; hold until ICR \\
  \hline
  63:5 & --   & --                                                              & Reserved \\
  \hline
\end{tabularx}
\end{table}

\paragraph{RIS} -:
Raw Interrupt Status Register. Read-only. Reflects latched interrupt conditions
independent of masking. Set by hardware edge events or by ISR write. Cleared by
ICR write. Writing has no effect (use ISR/ICR instead).
\asregkopf{Raw Interrupt Status Register}{$0000000000000000_H$}
\asregfourxsixteen{\multicolumn{16}{|c|}{0}}
                             {\multicolumn{16}{|c|}{r}}
                             {\multicolumn{16}{|c|}{0}}
                             {\multicolumn{16}{|c|}{r}}
                             {\multicolumn{16}{|c|}{0}}
                             {\multicolumn{16}{|c|}{r}}
                             {0&0&0&0&0&0&0&0&0&0&0&TI&ER&TH&HA&DO}
                             {r&r&r&r&r&r&r&r&r&r&r&rh&rh&rh&rh&rh}
\asregdesc{- & 63:5 & r & Reserved \\
           \hline
           TI & 4 & rh & TIMEOUT. Timeout threshold reached. \\
           \hline
           ER & 3 & rh & ERROR. Error latched in kernel. \\
           \hline
           TH & 2 & rh & TXHALF. TX FIFO below half-full threshold. \\
           \hline
           HA & 1 & rh & RXHALF. RX FIFO at or above half-full threshold. \\
           \hline
           DO & 0 & rh & DONE. Transfer completed. Latches on rising edge of \texttt{stat\_done\_s}.}

\paragraph{IMSC} -:
Interrupt Mask Control Register. Each bit enables (1) or disables (0) the
corresponding interrupt source. Reset value \texttt{0xFFFFFFFFFFFFFFFF} means
all interrupts are enabled after reset.
\asregkopf{Interrupt Mask Control Register}{$FFFFFFFFFFFFFFFF_H$}
\asregfourxsixteen{\multicolumn{16}{|c|}{IMSC[63:48]}}
                             {\multicolumn{16}{|c|}{rw}}
                             {\multicolumn{16}{|c|}{IMSC[47:32]}}
                             {\multicolumn{16}{|c|}{rw}}
                             {\multicolumn{16}{|c|}{IMSC[31:16]}}
                             {\multicolumn{16}{|c|}{rw}}
                             {0&0&0&0&0&0&0&0&0&0&0&TI&ER&TH&HA&DO}
                             {rw&rw&rw&rw&rw&rw&rw&rw&rw&rw&rw&rw&rw&rw&rw&rw}
\asregdesc{- & 63:5 & rw & Reserved (kept for forward compatibility; writes ignored by hardware) \\
           \hline
           TI & 4 & rw & TIMEOUT mask. 1 = interrupt enabled. \\
           \hline
           ER & 3 & rw & ERROR mask. \\
           \hline
           TH & 2 & rw & TXHALF mask. \\
           \hline
           HA & 1 & rw & RXHALF mask. \\
           \hline
           DO & 0 & rw & DONE mask.}

\paragraph{MIS} -:
Masked Interrupt Status Register. Read-only. Returns the effective interrupt
status as seen by the CPU. $\text{MIS} = \text{RIS} \;\&\; \text{IMSC}$.
\asregkopf{Masked Interrupt Status Register}{$0000000000000000_H$}
\asregfourxsixteen{\multicolumn{16}{|c|}{0}}
                             {\multicolumn{16}{|c|}{r}}
                             {\multicolumn{16}{|c|}{0}}
                             {\multicolumn{16}{|c|}{r}}
                             {\multicolumn{16}{|c|}{0}}
                             {\multicolumn{16}{|c|}{r}}
                             {0&0&0&0&0&0&0&0&0&0&0&TI&ER&TH&HA&DO}
                             {r&r&r&r&r&r&r&r&r&r&r&rh&rh&rh&rh&rh}
\asregdesc{- & 63:5 & r & Reserved \\
           \hline
           TI & 4 & rh & TIMEOUT masked \\
           \hline
           ER & 3 & rh & ERROR masked \\
           \hline
           TH & 2 & rh & TXHALF masked \\
           \hline
           HA & 1 & rh & RXHALF masked \\
           \hline
           DO & 0 & rh & DONE masked}

\paragraph{ICR} -:
Interrupt Clear Register. Write-only. Writing 1 to a bit clears the corresponding
RIS bit. Writing 0 has no effect. Reading returns 0. The interrupt should be
cleared only after the condition has been handled (FIFO drained / refilled)
to avoid re-triggering.
\asregkopf{Interrupt Clear Register}{$0000000000000000_H$}
\asregfourxsixteen{\multicolumn{16}{|c|}{0}}
                             {\multicolumn{16}{|c|}{r}}
                             {\multicolumn{16}{|c|}{0}}
                             {\multicolumn{16}{|c|}{r}}
                             {\multicolumn{16}{|c|}{0}}
                             {\multicolumn{16}{|c|}{r}}
                             {0&0&0&0&0&0&0&0&0&0&0&TI&ER&TH&HA&DO}
                             {r&r&r&r&r&r&r&r&r&r&r&w&w&w&w&w}
\asregdesc{- & 63:5 & r & Reserved \\
           \hline
           TI & 4 & w & Write 1 to clear RIS.TIMEOUT \\
           \hline
           ER & 3 & w & Write 1 to clear RIS.ERROR \\
           \hline
           TH & 2 & w & Write 1 to clear RIS.TXHALF \\
           \hline
           HA & 1 & w & Write 1 to clear RIS.RXHALF \\
           \hline
           DO & 0 & w & Write 1 to clear RIS.DONE}

\paragraph{ISR} -:
Interrupt Set Register. Write-only. Provided for software testing of the
interrupt path. Writing 1 forces the corresponding RIS bit to 1.
Writing 0 has no effect. Reading returns 0.
\asregkopf{Interrupt Set Register}{$0000000000000000_H$}
\asregfourxsixteen{\multicolumn{16}{|c|}{0}}
                             {\multicolumn{16}{|c|}{r}}
                             {\multicolumn{16}{|c|}{0}}
                             {\multicolumn{16}{|c|}{r}}
                             {\multicolumn{16}{|c|}{0}}
                             {\multicolumn{16}{|c|}{r}}
                             {0&0&0&0&0&0&0&0&0&0&0&TI&ER&TH&HA&DO}
                             {r&r&r&r&r&r&r&r&r&r&r&w&w&w&w&w}
\asregdesc{- & 63:5 & r & Reserved \\
           \hline
           TI & 4 & w & Write 1 to force-set RIS.TIMEOUT (test only) \\
           \hline
           ER & 3 & w & Write 1 to force-set RIS.ERROR (test only) \\
           \hline
           TH & 2 & w & Write 1 to force-set RIS.TXHALF (test only) \\
           \hline
           HA & 1 & w & Write 1 to force-set RIS.RXHALF (test only) \\
           \hline
           DO & 0 & w & Write 1 to force-set RIS.DONE (test only)}

% =============================================================================
\subsection{Software Programming Sequence}
% =============================================================================

\subsubsection{Normal Single Read (e.g.\ 0x0B Fast Read, 8\,B)}

\begin{enumerate}
  \item Write \texttt{CLKDIV} to set the desired SCK frequency.
  \item Write \texttt{CMD} = \texttt{0x0B}.
  \item Write \texttt{ADDR} = target flash address.
  \item Write \texttt{LEN} = \texttt{8}.
  \item Write \texttt{DUMMY} = \texttt{8}.
  \item Optionally write \texttt{IMSC} to unmask DONE interrupt (bit 0).
  \item Write \texttt{CTRL} = \texttt{0x08} (START = 1, all mode bits = 0 for single mode).
  \item Poll \texttt{RIS.DO = 1} or wait for \texttt{qspi\_irq\_o}.
  \item Write \texttt{ICR} = \texttt{0x01} to acknowledge DONE.
  \item Read \texttt{RXDATA} to retrieve received bytes (one 64-bit read per 8 bytes).
\end{enumerate}

\subsubsection{Quad Read (e.g.\ 0x6B Quad Output Fast Read, 8\,B)}

Same as above, with \texttt{CTRL} = \texttt{0x18} (QUAD=1, START=1, bit\,4 and bit\,3 set)
and \texttt{DUMMY} = \texttt{8}.

\subsubsection{TX Write (e.g.\ 0x02 Page Program, 8\,B)}

\begin{enumerate}
  \item Write \texttt{TXDATA} with up to 8 bytes (pushed to TX FIFO).
  \item Program \texttt{CMD}, \texttt{ADDR}, \texttt{LEN}, \texttt{DUMMY} as needed.
  \item Write \texttt{CTRL} = \texttt{0x08} (START = 1).
  \item Poll \texttt{RIS.DO} or wait for IRQ. Clear with ICR.
\end{enumerate}

\subsubsection{XIP Enable}

\begin{enumerate}
  \item Configure \texttt{CLKDIV}, \texttt{DUMMY}, \texttt{CMD} (e.g.\ 0xEB).
  \item Write \texttt{CTRL} = \texttt{0x50} (QUAD=1, XIP=1, START=0).
  \item XIP becomes active after the next DONE. \texttt{STATUS.XI} = 1.
  \item CPU reads to the flash address range are now automatically serviced.
  \item To exit XIP: write \texttt{CTRL.XIP = 0}.
\end{enumerate}

